% ===================================================================================================================
% Evaluation -> Methodology, in the structure of Phi (ISCA'25, Sec. 5.1): Models and Datasets / Implementation /
% Baselines, Table I = architecture setup, Table II = configuration and performance of all designs on VGG16.
% Items marked \tbd{...} wait for the pruned models (Rakshitha) and are printed in red.
% Areas in Table II: hw/same_node_table.py -> results/same_node_45nm.json (core RTL + CACTI SRAM, 45 nm).
% ===================================================================================================================
\section{Evaluation}
\label{sec:eval}

In this section, we first detail our evaluation methodology. Next, we justify the weight-stationary
dataflow over input- and output-stationary alternatives. We then evaluate the performance, energy and
area of TIDES when applied to SNN models. Additionally, we report the compression achieved by packing
and illustrate the reduction in memory traffic.

\subsection{Methodology}
\label{sec:method}

\textbf{Evaluation Models and Datasets.} We evaluate TIDES on spiking CNNs (VGG16 and ResNet18) on
CIFAR-10 with $T=8$ timesteps. We obtain the models from the open-source NeuroFlex
repository~\cite{neuroflex}, which converts INT8-quantized ANNs into PASCAL SNNs whose spiking execution
reproduces the quantized ANN exactly, and evaluate them using PyTorch; the dense VGG16 reaches 95.62\%
test accuracy. To expose weight sparsity, we prune the weights to \tbd{90\%} sparsity with
\tbd{magnitude pruning} and fine-tune the models, reaching \tbd{--} and \tbd{--} accuracy on VGG16 and
ResNet18, respectively. Then, we obtain the input spike trains and INT8 weights of every convolutional
and fully connected layer and utilize this information in our experiments. The first layer, which
receives the analog image, is excluded for all designs, and a layer that follows average pooling is
executed exactly on the four pre-pooling spikes of each pooled input. For sensitivity studies, we
additionally generate synthetic layers of the same shapes with independently set weight sparsity,
fraction of silent neurons and per-timestep firing probability.

\textbf{TIDES Implementation.} We developed a cycle-level simulator that takes the spike trains and
weights of the model as input, modeling the behavior of the packer, the PE array, the edge unit and the
LIF units to estimate the runtime performance of TIDES. We implement all components in Verilog; in a
closed-loop Verilator~\cite{verilator} testbench, the RTL reproduces the partial sums, output spikes
and cycle count of the simulator exactly on every layer of both networks at batch sizes 1 and 16. The
components are synthesized with Yosys and ABC~\cite{yosys} on the Nangate 45\,nm open cell
library~\cite{nangate45}, and OpenSTA~\cite{opensta} provides timing and static/dynamic power from
gate-level switching activity of real layers. We use CACTI 7.0~\cite{cacti7} to evaluate the area and
power of the on-chip buffers in a 45\,nm process, with access counts from the simulator. We use the DRAM
energy per bit derived with DRAMsim3~\cite{dramsim3} in~\cite{prosperity2025} to estimate the off-chip
DRAM energy. Tbl.~\ref{tab:setup} shows the setup of the TIDES architecture.

\begin{table}[t]
  \centering
  \caption{TIDES Architecture Setup.}
  \label{tab:setup}
  \small
  \setlength{\tabcolsep}{4pt}
  \begin{tabular}{@{}ll@{}}
    \toprule
    Parameter & Value \\
    \midrule
    Tile Size           & $R=4$ inputs; output 64 filters $\times$ 128 pixels \\
    Packing             & $F=4$ filters/column; $G=4$ positions/group \\
    Computation Array   & $4\times4$ PEs $\times$ 8 adders; 16 LIF units \\
    Edge Unit           & 4 FIFOs $\times$ 80 entries; 4 RMW/column/cycle \\
    On-chip Buffer Size & 192\,KB Output; 64\,KB Staging; 4\,KB Group \\
    DDR4 DRAM           & 4 Channels, 64\,GB/s \\
    \bottomrule
  \end{tabular}
\end{table}

\begin{table*}[t]
  \centering
  \caption{Comparison of TIDES with Baselines on VGG16 (CIFAR-10). Area includes the on-chip buffers.}
  \label{tab:baselines}
  \small
  \begin{tabular}{@{}lcccccc@{}}
    \toprule
                 & Freq.  & Tech. & Area    & Throughput & Energy Efficiency & Area Efficiency \\
                 & (MHz)  & (nm)  & (mm$^2$) & (GOP/s)   & (GOP/J)           & (GOP/s/mm$^2$) \\
    \midrule
    PTB~\cite{lee2022ptb}              & 500 & 45 & 1.232 & \tbd{--} (1.00$\times$) & \tbd{--} (1.00$\times$) & \tbd{--} (1.00$\times$) \\
    SpikeX~\cite{xu2025spikex}         & 500 & 45 & 0.747 & \tbd{--} & \tbd{--} & \tbd{--} \\
    SATO~\cite{liu2022sato}            & 500 & 45 & 0.354 & \tbd{--} & \tbd{--} & \tbd{--} \\
    SpikeFlow~\cite{wang2025spikeflow} & 500 & 45 & 0.207 & \tbd{--} & \tbd{--} & \tbd{--} \\
    LoAS~\cite{yin2024loas}            & 500 & 45 & 1.478 & \tbd{--} & \tbd{--} & \tbd{--} \\
    Prosperity~\cite{prosperity2025}   & 500 & 45 & 1.035 & \tbd{--} & \tbd{--} & \tbd{--} \\
    TIDES                              & 500 & 45 & 2.128 & \tbd{--} & \tbd{--} & \tbd{--} \\
    \bottomrule
  \end{tabular}
\end{table*}

\textbf{Baselines.} We compare with the SNN accelerators PTB~\cite{lee2022ptb},
SpikeX~\cite{xu2025spikex}, SATO~\cite{liu2022sato}, SpikeFlow~\cite{wang2025spikeflow},
LoAS~\cite{yin2024loas} and Prosperity~\cite{prosperity2025}, which we discussed in
Sec.~\ref{sec:background}. We compare their configuration and performance on VGG16 with CIFAR-10 in
Tbl.~\ref{tab:baselines}. For a fair comparison, we keep the frequency and technology the same for all
baselines and TIDES: the core of each baseline is implemented in RTL and taken through the same
synthesis and power flow, and its on-chip buffers, sized as in its paper, are evaluated with CACTI 7.0
in the same process. We define operations (OP) consistently across all baselines and TIDES, where each
OP corresponds to an accumulation for a `1' element in the spike bit matrix. We simulate PTB, SATO and
Prosperity with the official simulators released with Prosperity~\cite{prosperity2025}, and SpikeX,
SpikeFlow and LoAS in our framework. On our VGG16 traces, the official simulators give Prosperity a
9.6$\times$ speedup over PTB, close to the 9.4$\times$ reported in~\cite{prosperity2025}.
